\par\noindent\textbf{(1)} 本征矩阵H=[ω₁,g;g,ω₂]，令平均频率 ω̄=(ω₁+ω₂)/2 与失谐 Δ=ω₁−ω₂，其本征频率 λ\_±=ω̄±√[(Δ/2)²+g²]。当调谐Δ=0，两模式最小频率间距2g，避免本征频率交叉；g→0才恢复两独立模。
\par\noindent\textbf{(2)} 去掉公共相位后演化为 H′=(Δ/2)σ\_z+gσ\_x，故 exp(−iH′t)=cos(Ω\_R t/2)I−2iH′sin(Ω\_Rt/2)/Ω\_R，其中 Ω\_R=√(Δ²+4g²)。初态(1,0)时，另一腔幅度 −i(2g/Ω\_R)sin(Ω\_Rt/2)，能量占比 P₂(t)=4g² sin²(Ω\_Rt/2)/(Δ²+4g²)。
\par\noindent\textbf{(3)} 最大交换概率 P₂,max=4g²/(Δ²+4g²)，只在精确共振时达到1；振荡周期为 2π/Ω\_R。大失谐 |Δ|≫g 时，P₂,max≈4g²/Δ²，能量主要留在原腔；可通过谱中避交叉间隙和拍振周期的联合测量反推g。
\par\noindent\textbf{(4)} 若两腔具有相同的振幅衰减率κ，则H\_eff=H−iκI，两本征频率均向复平面下移−iκ，避交叉实频隙不变。总归一化模能量为 e\textasciicircum{}\{−2κt\}，归一化条件交换形状不变；当两腔损耗不等时非Hermitian特征根不同且可能出现exceptional point，不能沿用等损耗式。
